\documentclass{article}
\usepackage[main]{neurips_2026}
\usepackage[utf8]{inputenc}
\usepackage[T1]{fontenc}
\usepackage{amsmath,amsfonts,amssymb}
\usepackage[hidelinks]{hyperref}
\usepackage{url,booktabs,microtype,graphicx,xcolor}
\renewcommand{\ttdefault}{pcr}
\title{Single-run SHAP reliability under finite computation}
\author{}
\begin{document}
\maketitle
\begin{abstract}
A stochastic SHAP attribution is a numerical estimate, even when displayed as a list of reasons for a prediction. We assess when a single run preserves a fixed empirical-game attribution under finite CPU budgets. Across three public tabular datasets, three fixed classifiers and two estimators, we compare 97,200 explanations with 540 exactly enumerated references. Our declared adequacy event requires normalized L1 error at most .20 and recovery of at least two of three signed salient features. Lowest-budget Kernel success ranges from .7794 to 1 across fitted cells; every highest-budget condition has observed success 1, with a conservative boundary interval [.8247,1] over 60 input clusters. This does not certify population or individual reliability. Requiring all three signed features reduces highest-budget success as low as .9289, and a lower-budget input has 0/30 successes despite high panel-average performance. Repeatability and reference accuracy are separately measured. The main experiment uses 7.493 CPU hours; interrupted elapsed clocks are retained and cannot establish deployment latency. The results support conditional numerical adequacy for the tested games, budgets and readings, with explicit failure cases and limits on transfer to consequential decisions.
\end{abstract}
\input{paper/sections/introduction.tex}
\input{paper/sections/related_work.tex}
\input{paper/sections/attribution_target.tex}
\input{paper/sections/assessment.tex}
\input{paper/sections/study_design.tex}
\input{paper/sections/results.tex}
\input{paper/sections/discussion.tex}
\input{paper/sections/scope.tex}
\clearpage
\bibliographystyle{plainnat}
\bibliography{paper/references}
\clearpage
\appendix
\input{paper/sections/pilot_accounting.tex}
\input{paper/sections/reproduction.tex}
\clearpage
\input{paper/sections/technical_definitions.tex}
\clearpage
\input{paper/sections/implementation_checks.tex}
\clearpage
\input{paper/sections/background.tex}
\clearpage
\input{paper/sections/numerical_appendix.tex}
\clearpage
\input{paper/checklist.tex}
\end{document}
